% Bounded continuation: preserved legacy results, age-only recombination, and development diagnostics.
The following pilot and age-only recombination describe the historical audit stage. The completed twenty-quarter Freddie rebuild, new scores, bridge and uncertainty diagnostics are reported in Appendix~\ref{app:correcteddiagnostics}. The ingestion repairs established two code defects without, at that stage, establishing their full empirical effect. The prior-balance calculation now orders each loan's observations by reporting month; an unsorted synthetic history reproduced the former error. Quarter-cell mean ages are now combined using loan counts rather than equal weights for quarter means. At that historical stage, checks for duplicate loan-months, gaps, first-observation payoff fallback and overlapping file membership were incomplete across the full panel. Chronological ordering alone does not establish that a lag is the immediately preceding monthly balance. The subsequent Freddie rebuild and dependent reruns are complete; original-principal fallback remains an explicit assumption where the preceding balance is unobserved.

A bounded raw-source pilot checks the 2017Q1 Freddie acquisition pair through September 2025. The 32-field schema and source-byte checks pass for 290,601 origination records and 16,516,599 performance records. Under the existing fixed-rate eligibility rules, 290,558 loans contribute 16,513,080 loan-months and 4,847 cohort-months. Both the legacy and repaired production calculations match an independent chronological reference within the frozen numerical tolerance for balances, original principal, payoff principal, counts, and mean age. This quarter therefore shows no measured payoff change from the ordering repair. Its eligible histories have no duplicate loan-months, backward ordering, missing ages, or invalid balances, but contain 63 missing calendar months. No payoff has a nonadjacent observed prior balance. For 139 first-observed payoffs, both implementations instead substitute original principal, totaling \$33.659 million; agreement with that fallback does not verify the unobserved preceding balance. This single-quarter result does not establish the other nineteen quarters' integrity or reconcile cross-file overlap, and its eligibility includes terms other than 30 years.

The saved and fresh simulation samples each contain 75,000 individual loan ages, with no missing ages and 65 distinct values from 0 through 64 months. Their simulations do not substitute the repaired pooled mean ages for these individual inputs. Both the saved \$42.6084 billion pair and the fresh \$42.1761 billion pair nevertheless use the legacy panel-derived 4.991\% annual floor. The cached transition matrix is another panel-related input. The legacy floor annualizes payoff principal divided by contemporaneous end exposure; predictive scoring instead uses lagged opening exposure, and simulated hazards apply to active beginning balances. An audited opening-risk-set calculation and coverage reconciliation are required before treating a rebuilt floor as an observed monthly payoff probability. None of the fixed-floor replays establishes a corrected headline estimate or the direction of a future recalibration.

The completed Fannie exercise is an age-only recombination of 24 preserved acquisition-quarter files spanning 2017Q1--2022Q4. The legacy rule reproduces all 23,443 rows of the saved 2017--2021-origination panel, and 28 saved Freddie/Fannie floor-grid entries pass replay checks. Replacing only the age-combination rule conserves keys, counts, exposure, original principal, payoff principal, monthly payoff rates, burnout, and modal state. It changes 20,906 pooled mean ages. Count weighting recovers the pooled known-age mean only when the reported count is the count underlying each finite quarter mean; individual missing-age counts are unavailable in these aggregates. Entirely missing quarter means are excluded from the age average without removing their financial amounts. A cohort mean-age threshold does not restrict every constituent loan's age.

For the recombined Fannie 2017--2021-origin panel pooled over calendar 2017--2019, at a note-minus-market gap of at most $-25$ bp, the descriptive annualized payoff rate changes from 5.522\% to 5.544\% at mean age at least twelve months, with 421 versus 489 qualifying cohort-months. At mean age at least 24 months, it changes from 10.136\% to 9.523\%, with 71 versus 92 cohort-months. These rates retain saved payoff numerators and the end-balance denominator; they are not raw-source-corrected payoff estimates. The pooled 2017--2019 object is distinct from the 2018-only Freddie anchor and the older-origination Fannie exercise below.

That separate exercise selects origins through 2016 observed during calendar 2018 in the later Fannie acquisition files. At mean age at least 24 months and gap at most $-25$ bp, count weighting reclassifies 99.05\% of previously qualifying exposure below the age threshold; it does not delete those balances. Remaining support is 183 cohort-months in 23 strata, totaling 0.443588 billion dollar-months of end exposure. Concentration corresponds to 5.10 equally weighted strata, not 5.10 independent observations. There is no support at a $-200$ bp gap. The descriptive 5.5621\% annualized payoff rate describes a thin, selected later-acquisition tail. It does not establish a representative mature-loan floor or a sale-related moving share, and it does not validate the distinct 3.5\% prediction candidate.

The separate development exercise uses the preserved Freddie panel and the original supported 239 cohorts. Its fixed grid contains annual floors of 0, 2, 3.5, and 4.991\% and PSA speeds of 50, 75, and 100. Selection uses January--December 2023 only, initialized in December 2022; both permitted selection families choose 3.5\% and 100 PSA. On the already-examined January 2024--September 2025 evaluation support, observed voluntary payoff principal is \$131.3033 billion. The original response-on model predicts \$161.7842 billion, the development candidate \$128.8461 billion, and the simple 2023 cohort-speed comparator \$121.4362 billion. Their exposure-weighted monthly-rate RMSEs are respectively 0.13389, 0.11014, and 0.08679 percentage points. The candidate reduces the original error by 17.7376\% but still loses to the comparator. Its response-off prediction is \$160.7063 billion with RMSE 0.15355, versus \$172.6429 billion and 0.15700 for the original response-off model. These are legacy-input conditional scores, not new SOMA simulation results.

On the original response-on leg, the 4.991\% floor binds on 91.5353\% of evaluation exposure and alone implies \$160.4220 billion of payoff principal, explaining most of the original aggregate overprediction mechanically. Predictions condition on realized opening balances and mortgage rates; burnout advances using predicted payoffs. The approximate internal balance recursion omits scheduled amortization and defaults in both original and candidate prediction checks. Selection on 2023 does not make the already-seen 2024--September 2025 outcomes an untouched holdout. The candidate is retained for future testing without promotion into the paired results. A future evaluation requires frozen sources, parser, sample, parameters, and treatment rules before genuinely new outcomes are accessed. Better observed fit would still not validate the response-off counterfactual, imported elasticity, sale share, or full-SOMA transport.
